%% ============================================================
%% Bachelor's Thesis Proposal — Institute for International Business (IIB), WU Wien
%% Arthur Wienerroither — Student ID: h12322363
%% Target: 3–4 pages text + 1 page references (10 top-ranked sources)
%% [SOURCE: ...] markers = slots for articles from the Scopus search.
%% ============================================================
\documentclass[12pt, a4paper]{article}

\usepackage[utf8]{inputenc}
\usepackage[T1]{fontenc}
\usepackage[english]{babel}
\usepackage{csquotes}
\usepackage[a4paper, margin=2.5cm]{geometry}
\usepackage{setspace}
\onehalfspacing
\usepackage[protrusion=true,expansion=false]{microtype}
\usepackage{booktabs}
\usepackage{array}
\usepackage[hidelinks]{hyperref}
\usepackage{xurl}
\usepackage{xcolor}

% Visible placeholder for sources to be pulled from Scopus
\newcommand{\src}[1]{\textcolor{red}{\textbf{[SOURCE: #1]}}}

\begin{document}

%% --- Header (compact — page budget is tight) ---
\begin{center}
  {\LARGE\bfseries When Do AI Investments Create Competitive Advantage?\par}
  \vspace{0.3cm}
  {\large A Systematic Review of Firm-Level Empirical Evidence (2015--2026)\par}
  \vspace{0.6cm}
  {\normalsize Bachelor's Thesis Proposal\par}
  {\normalsize Institute for International Business, WU Wien\par}
  {\normalsize Supervisor: PD Dr.\ Viktor Fredrich\par}
  \vspace{0.3cm}
  {\normalsize Arthur Wienerroither (h12322363) --- \today\par}
\end{center}
\vspace{0.5cm}

\section{Introduction and Relevance}

Artificial intelligence (AI), understood as machines that perform ``cognitive
functions that are usually associated with human minds, such as learning,
interacting, and problem solving'' (Raisch \& Krakowski, 2021, p.~192)\footnote{[Raisch \& Krakowski (2021), p.~192] ``Artificial intelligence (AI) refers to machines performing cognitive functions that are usually associated with human minds, such as learning, interacting, and problem solving.''},
has moved from research laboratories to the core of corporate investment agendas.
Firms across industries are committing substantial resources to AI. The
returns on these investments, however, remain uncertain. Nine of ten executives see a major
business opportunity in AI, yet ``most companies have a hard time
generating value with AI'' (Ransbotham et al., 2019, p.~1, as cited in Kemp, 2024,
p.~619)\footnote{[Kemp (2024), p.~619] ``In a recent survey by Ransbotham, Khodabandeh, Fehling, LaFountain, \& Kiron (2019: 1), nine out of 10 top managers reported that AI represents a large business opportunity for their firms and 43\% of executives reported having implemented AI in their organizations. The report also noted, however, that `most companies have a hard time generating value with AI.'''}.
Stock markets have on average even reacted negatively to firms' AI investments
(Lui, Lee, \& Ngai, 2022, as cited in Kemp, 2024, p.~619)\footnote{[Kemp (2024), p.~619] ``Lui, Lee, and Ngai (2022) offered empirical support for this concern, finding that markets penalize AI adoption at the firm level.''}.

Theory can make sense of this disconnect. AI ``is regarded as a general-purpose technology akin
to electricity, the steam engine, or the internet'' (Kemp, 2024, p.~620)\footnote{[Kemp (2024), p.~620] ``AI is regarded as a general-purpose technology akin to electricity, the steam engine, or the internet (Brynjolfsson \& Mitchell, 2017; Frey \& Osborne, 2013; Lynch, 2017).''},
and from the resource-based view a technology that every competitor can buy
should not, by itself, be a source of sustained competitive advantage (Barney,
1991, p.~110)\footnote{[Barney (1991), p.~110] ``If one firm can purchase these physical tools of production and thereby implement some strategies, then other firms should also be able to purchase these physical tools, and thus such tools should not be a source of sustained competitive advantage.''}. The advantage goes to the firm that combines the
technology with resources others cannot copy, such as its social relations
and culture (Barney, 1991, p.~110)\footnote{[Barney (1991), p.~110] ``Several firms may all possess the same physical technology, but only one of these firms may possess the social relations, culture, traditions, etc.\ to fully exploit this technology in implementing strategies.''}. This thesis asks under what conditions firms
achieve this combination, a question that matters to strategic management
theory and to practitioners deciding which AI budgets create value.

\section{Problem Definition and Research Gap}

Empirical research linking firm-level AI investments to performance outcomes
has grown quickly, but it is scattered across several literatures and points
in different directions. Babina, Fedyk, He, and
Hodson (2024, p.~1)\footnote{[Babina et al. (2024), p.~1] ``AI-investing firms experience higher growth in sales, employment, and market valuations. This growth comes primarily through increased product innovation.''} find higher growth and more product innovation among
AI-investing firms, concentrated among firms that were already large (Babina
et al., 2024, p.~3)\footnote{[Babina et al. (2024), p.~3] ``the positive relationship between AI investments and firm growth is much stronger among ex-ante larger firms''}. Brynjolfsson, Rock, and Syverson (2021, p.~334)\footnote{[Brynjolfsson, Rock \& Syverson (2021), p.~334] ``As firms adopt a new GPT, total factor productivity growth will initially be underestimated because capital and labor are used to accumulate unmeasured intangible capital stocks.''}
explain why the payoff often arrives late or not at all: before AI can raise
productivity, a firm must build assets that no balance sheet records, such as
retrained staff and reorganized workflows. Lui, Lee, and Ngai (2022, p.~373)\footnote{[Lui, Lee \& Ngai (2022), p.~373] ``The stock prices of the firms decrease by 1.77\% on the day of the announcement.''}
find share prices dropping by 1.77\% on average when firms announce AI
investments. The drop is larger where the firm's IT capability is weak (Lui
et al., 2022, p.~373)\footnote{[Lui, Lee \& Ngai (2022), p.~373] ``Nonmanufacturing firms and firms with weak information technology capabilities or low credit ratings suffer a more negative impact compared with other firms.''}.

Recent theory explains why the results differ: the payoff depends on
conditions inside the firm. A firm gains an advantage its rivals cannot copy
when it \emph{situates} AI in its own experience, structure, and
relationships (Kemp,
2024, p.~618)\footnote{[Kemp (2024), p.~618] ``The paper highlights the organizational activities involved in situating AI --- specifically, grounding, bounding, and recasting.''}. Outcomes also
depend on how a firm balances automation, where AI replaces human work, and
augmentation, where AI supports it (Raisch
\& Krakowski, 2021, p.~192)\footnote{[Raisch \& Krakowski (2021), p.~192] ``Overemphasizing either augmentation or automation fuels reinforcing cycles with negative organizational and societal outcomes.''}. Early empirical work supports this view:
machine learning beat an older technology only when its users knew the domain
well (Choudhury, Starr, \& Agarwal, 2020, p.~1382)\footnote{[Choudhury, Starr \& Agarwal (2020), p.~1382] ``Domain expertise complements ML by correcting for the (strategic) incompleteness of the input to the ML tool, while vintage-specific skills ensure the ability to properly operate the technology.''}.

No systematic review has consolidated this evidence. A prominent review
asks how firms adopt and use AI, not under which conditions AI leads to
strategic outcomes (Enholm, Papagiannidis, Mikalef, \& Krogstie, 2022,
p.~1709)\footnote{[Enholm et al. (2022), p.~1709] ``This study provides a systematic literature review that attempts to explain how organizations can leverage AI technologies in their operations and elucidate the value-generating mechanisms. Our analysis synthesizes the current literature and highlights: (1) the key enablers and inhibitors of AI adoption and use; (2) the typologies of AI use in the organizational setting; and (3) the first- and second-order effects of AI.''}. The same review concedes that ``it is still not
clear if and how AI can help organizations achieve financial performance
gains'' (Enholm et al., 2022, p.~1728)\footnote{[Enholm et al. (2022), p.~1728] ``Nevertheless, there is a long chain of causal associations, and to date it is still not clear if and how AI can help organizations achieve financial performance gains.''}. This thesis aims to close that gap. It collects and
compares the firm-level evidence on the conditions under which AI
investments turn into competitive advantage. These conditions include
organizational capabilities, complementary assets, the design of the
investment itself, and the firm's environment.

\section{Research Question and Objectives}

The thesis will answer the following research question:

\begin{quote}
\emph{Under what conditions do firm-level AI investments translate into
competitive advantage and superior firm performance? --- A systematic review of
empirical evidence (2015--2026).}
\end{quote}

Three filters narrow the question. First, the study must appear between 2015
and 2026. Second, it must study whole firms, not individuals or single tasks.
Third, it must come from a top-ranked peer-reviewed journal (AJG~4* or~4, or
VHB~A+ or~A). A paper from a lower-ranked journal is included only if it offers
something the thesis needs, such as a unique data set. Every exception is
documented with its reason. The thesis then has three goals: to map what the
evidence says about how AI investment affects firm performance; to record the
conditions that make this effect positive, negative, or absent; and to compare
that evidence with current theory to find directions for future research.

\section{Method}

The thesis employs a systematic literature review (SLR) designed for maximal
transparency and reproducibility (Tranfield, Denyer, \& Smart, 2003,
p.~209)\footnote{[Tranfield, Denyer \& Smart (2003), p.~209] ``Systematic reviews differ from traditional `narrative' reviews by adopting a replicable, scientific and transparent process [...] that aims to minimize bias through exhaustive literature searches of published and unpublished studies and by providing an audit trail of the reviewers decisions, procedures and conclusions.''}.
The search will be conducted in Scopus (WU license); pilot searches via the
Scopus API calibrated the following query:

\begin{quote}\small\raggedright
\texttt{TITLE(("artificial intelligence" OR "machine learning" OR "deep
learning" OR "AI")) AND TITLE-ABS-KEY(("invest" OR "invests" OR "invested" OR
"investing" OR "investment" OR "investments" OR "investor" OR "investors" OR
adopt* OR implement*) AND ("competitive advantage" OR "firm performance" OR
"financial performance" OR "firm productivity" OR "firm value" OR "market
value" OR "firm growth")) AND SUBJAREA(BUSI OR ECON OR DECI) AND PUBYEAR > 2014
AND PUBYEAR < 2027 AND DOCTYPE(ar) AND LANGUAGE(english) AND SRCTYPE(j)}
\end{quote}

\noindent As of 6~July~2026, the query returns 425 records, which lies inside
the required corridor of 50--500 empirical studies. Before the search, three
empirical studies were set as test cases that any sound query must find:
Babina et al. (2024), Lui et al. (2022), and Krakowski, Luger, and Raisch
(2023). The query finds all three, so it does not miss known relevant work. Included are empirical, firm-level, peer-reviewed journal
articles in English published 2015--2026; excluded are conceptual papers,
individual- or task-level studies, and non-ranked outlets. The screening
process will be documented in a PRISMA-style flow diagram, so every decision
can be reproduced.

Each included study is then coded into a state-of-the-art table: authors and
year, journal and ranking, theoretical lens, method and sample, AI investment
measure, performance outcome measure, direction of the effect, and identified
conditions. This table is the core of the thesis.

\section{Preliminary Structure of the Thesis}

\begin{center}
\small
\begin{tabular}{@{}>{\raggedright\arraybackslash}p{0.29\textwidth}>{\raggedright\arraybackslash}p{0.63\textwidth}@{}}
\toprule
\textbf{Chapter} & \textbf{Content} \\
\midrule
1.~Introduction & Relevance of the topic, the research gap, and the research question. \\
2.~Theoretical Background & AI as a general-purpose technology, resource-based reasoning, situated AI, and the automation--augmentation paradox. \\
3.~Method & SLR protocol, search strategy, and coding scheme. \\
4.~Results & Descriptive overview of the corpus and the state-of-the-art table. \\
5.~Discussion & A framework of the conditions for AI-driven competitive advantage, compared with current theory. \\
6.~Limitations & Constraints of the search, the corpus, and the reviewed studies. \\
7.~Conclusion and Future Outlook & Summary of the findings and directions for future research. \\
\bottomrule
\end{tabular}
\end{center}

%% ============================================================
\newpage
\section*{References}

\begingroup\small
\setlength{\parskip}{0pt}
\begin{itemize}\setlength{\itemsep}{2pt}
  \item Babina, T., Fedyk, A., He, A., \& Hodson, J. (2024). Artificial
  intelligence, firm growth, and product innovation. \emph{Journal of Financial
  Economics, 151}, 103745. \url{https://doi.org/10.1016/j.jfineco.2023.103745}
  \hfill [AJG 4*]
  \item Barney, J. (1991). Firm resources and sustained competitive advantage.
  \emph{Journal of Management, 17}(1), 99--120.
  \url{https://doi.org/10.1177/014920639101700108} \hfill [AJG 4*]
  \item Brynjolfsson, E., Rock, D., \& Syverson, C. (2021). The productivity
  J-curve: How intangibles complement general purpose technologies.
  \emph{American Economic Journal: Macroeconomics, 13}(1), 333--372.
  \url{https://doi.org/10.1257/mac.20180386} \hfill [AJG 4]
  \item Choudhury, P., Starr, E., \& Agarwal, R. (2020). Machine learning and
  human capital complementarities: Experimental evidence on bias mitigation.
  \emph{Strategic Management Journal, 41}(8), 1381--1411.
  \url{https://doi.org/10.1002/smj.3152} \hfill [AJG 4*]
  \item Enholm, I. M., Papagiannidis, E., Mikalef, P., \& Krogstie, J. (2022).
  Artificial intelligence and business value: A literature review.
  \emph{Information Systems Frontiers, 24}(5), 1709--1734.
  \url{https://doi.org/10.1007/s10796-021-10186-w} \hfill [AJG 3]
  \item Kemp, A. (2024). Competitive advantage through artificial intelligence:
  Toward a theory of situated AI. \emph{Academy of Management Review, 49}(3),
  618--635. \url{https://doi.org/10.5465/amr.2020.0205} \hfill [AJG 4*]
  \item Krakowski, S., Luger, J., \& Raisch, S. (2023). Artificial
  intelligence and the changing sources of competitive advantage.
  \emph{Strategic Management Journal, 44}(6), 1425--1452.
  \url{https://doi.org/10.1002/smj.3387} \hfill [AJG 4*]
  \item Lui, A. K. H., Lee, M. C. M., \& Ngai, E. W. T. (2022). Impact of
  artificial intelligence investment on firm value. \emph{Annals of Operations
  Research, 308}(1--2), 373--388. \url{https://doi.org/10.1007/s10479-020-03862-8}
  \hfill [AJG 3]
  \item Raisch, S., \& Krakowski, S. (2021). Artificial intelligence and
  management: The automation--augmentation paradox. \emph{Academy of Management
  Review, 46}(1), 192--210. \url{https://doi.org/10.5465/amr.2018.0072} \hfill [AJG 4*]
  \item Tranfield, D., Denyer, D., \& Smart, P. (2003). Towards a methodology
  for developing evidence-informed management knowledge by means of systematic
  review. \emph{British Journal of Management, 14}(3), 207--222.
  \url{https://doi.org/10.1111/1467-8551.00375} \hfill [AJG 4]
\end{itemize}
\endgroup

\end{document}
